\documentclass[../../main.tex]{subfiles}
\begin{document}

\begin{flushright}
\footnotesize\textit{【自動文字起こし・要確認】}
\end{flushright}

[解] $0 < a \le 2$ $\cdots$ ①, $t = \dfrac{1}{a}$ とおくと ①から $\dfrac{1}{2} \le t$ $\cdots$ ② である。

(1) 題意がみたされるのは、
\[ \sqrt{x} = 2tx + 1 - t \quad \cdots \text{③} \]
が $x$ について実数解を持つこと。$x = \dfrac{1}{2}$ は（$\sqrt{1/2} = 1$ となり不適だから、）以下 $x \neq 1/2$ とする。③より
\[ \dfrac{\sqrt{x}-1}{2x-1} = t \quad \cdots \text{④} \]
④の左辺を $f(x)$ とおく。
\[ f'(x) = \dfrac{\dfrac{1}{2\sqrt{x}}(2x-1) - 2(\sqrt{x}-1)}{(2x-1)^2} = \dfrac{-x + 2\sqrt{x} - 1/2}{\sqrt{x}(2x-1)^2} \]
だから、$p = \dfrac{2-\sqrt{2}}{2}, \ q = \dfrac{2+\sqrt{2}}{2}$ とおくと、
\[ f'(x) = \dfrac{-(\sqrt{x}-\sqrt{p})(\sqrt{x}-\sqrt{q})}{\sqrt{x}(2x-1)^2} \]
となるので、下表をうる。

\begin{center}
\begin{tabular}{|c|c|c|c|c|c|c|c|c|}
\hline
$x$ & $0$ & $\cdots$ & $p$ & $\cdots$ & $1/2$ & $\cdots$ & $q$ & $\cdots$ \\ \hline
$f'(x)$ & & $-$ & $0$ & $+$ & $/$ & $+$ & $0$ & $-$ \\ \hline
$f(x)$ & $1$ & $\searrow$ & $\dfrac{2+\sqrt{2}}{4}$ & $\nearrow$ & $/$ & $\nearrow$ & $\dfrac{2-\sqrt{2}}{4}$ & $\searrow$ \\ \hline
\end{tabular}
\end{center}

これと、$f(x) \to \pm\infty$ ($x \to 1/2 \mp 0$) (複号同順)、$f(x) \to 0$ ($x \to +\infty$) から、グラフの概形は右上図。これと $\dfrac{2-\sqrt{2}}{4} < \dfrac{1}{2} < \dfrac{2+\sqrt{2}}{4}$ 及び②から、④が $x$ について実数解を持つ条件は、
\[ \dfrac{2+\sqrt{2}}{4} \le t \iff a \le \dfrac{4}{2+\sqrt{2}} = 2(2-\sqrt{2}) \]

(2) (与式) $\iff 4x^2 + \{4(a-1)-a^2\}x + (a-1)^2 = 0 \iff 4x^2 - (a-2)^2 x + (a-1)^2 = 0$ $\cdots$ ⑤

この判別式を $D$ として、
\[ D = (a-2)^4 - 16(a-1)^2 = \{(a-2)^2 + 4(a-1)\}\{(a-2)^2 - 4(a-1)\} = a^2(a^2-8a+8) \quad \cdots \text{⑥} \]
だから、$t = 4-2\sqrt{2}$ として、①とから、
\begin{align}
&0 < a < t \text{ の時、} D \ge 0 \text{ から、⑤は実数解を持つ。} \\
&t < a \le 2 \text{ の時、} D < 0 \text{ から、⑤は実数解を持たない。} \quad \cdots \text{⑦}
\end{align}
となる。⑤の左辺を $g(x)$ とおく。

$1^\circ$ $0 < a < t$ の時

⑤の2解は、$g(1/2) \ge 0$、軸 $x = \dfrac{(a-2)^2}{8} \ge 0$ から、いずれも非負だから、$|\alpha| \le |\beta| \therefore \alpha \le \beta$ で、
\[ |\beta| = \beta = \dfrac{1}{8}\{(a-2)^2 + \sqrt{D}\} \quad \cdots \text{⑧} \]
\begin{align}
8\dfrac{d\beta}{da} &= 2(a-2) + \dfrac{4a(a^2-6a+4)}{2\sqrt{D}} \\
&= 2(a-2) + \dfrac{2(a^2-6a+4)}{\sqrt{a^2-8a+8}} \quad \cdots \text{⑨}
\end{align}
だから、
\[ \dfrac{d\beta}{da} \ge 0 \iff \dfrac{a^2-6a+4}{\sqrt{a^2-8a+8}} \ge 2-a \]
$0 < a \le 3-\sqrt{5}$ の時は 両辺0以上、$3-\sqrt{5} < a \le t$ の時は 左辺が負、右辺が正となり 不適だから、
$0 < a \le 3-\sqrt{5}$ のもとで2乗して、
\begin{align}
&(a^2-6a+4)^2 \ge (a-2)^2(a^2-8a+8) \\
&\iff a^4 - 12a^3 + 44a^2 - 48a + 16 \ge a^4 - 12a^3 + 44a^2 - 64a + 32 \\
&\iff a \ge 1
\end{align}
だが、$1 > 3-\sqrt{5}$ より、$\dfrac{d\beta}{da} \ge 0$ をみたす $a$ はなく、区間内で $\dfrac{d\beta}{da} < 0$。つまり $\beta$ は単調減少。
さらに、$\beta$ は明らかに連続だから、$0 < a \le t$ とすると、$\beta$ は $a=t$ の時
\[ \min |\beta| = \dfrac{1}{2}(3-2\sqrt{2}) \]
をとる。

$2^\circ$ $t \le a \le 2$ の時

$|\beta| = \dfrac{1}{8}\sqrt{(a-2)^4 - D} = \dfrac{1}{2}|a-1|$ で、$1 < t$ から、$a=t$ で $\min |\beta| = \dfrac{1}{2}(3-2\sqrt{2})$ をとる。

以上 $1^\circ, 2^\circ$ から、求める $\min |\beta| = \dfrac{1}{2}(3-2\sqrt{2})$ である。

\end{document}
